\section{Source Code and Reproducibility}

Supplementary files are maintained in the submitted project repository rather than duplicated in full here. The public repository is \url{https://github.com/Shui-Zhou/final-project-upload}; the assessed source snapshot is the submitted code bundle. This appendix indexes the source code, selected listings, data sources, figures, evaluation records, and dashboard operating guide. The listed entry points rebuild the processed data, run the dashboard locally, and trace the report figures and evaluation materials.

\subsection{Source-Code Index}

\begin{center}
\small
\begin{tabularx}{\textwidth}{p{4.3cm}X}
\toprule
Path & Purpose \\
\midrule
\path{src/data/build_canonical.py} & Builds the canonical incident parquet and its dictionary, provenance, and data-quality memo. \\
\path{src/data/build_borough_summary.py} & Aggregates the canonical incident table to the cells defined by borough, month, and incident group used by the dashboard. \\
\path{src/data/build_mobilisations.py} & Builds the mobilisation table and persists the \path{matches_canonical_incident} flag used for matched-only deployment summaries. \\
\path{src/data/build_station_locations.py} & Builds the incident-derived assignment-footprint station centroids. \\
\path{src/data/build_published_station_proximity.py} & Builds the public station-address postcode-geocoding recovery, D1 comparison table, borough proximity table, and station-proximity figures. \\
\path{src/data/build_report_evaluation_figures.py} & Builds report-side early-context, results, and robustness figures with JSON/markdown sidecar memos. \\
\path{src/data/build_closure_cohort_context.py} & Builds the Taylor closure-cohort descriptive cross-check figure, parquet, provenance JSON, and DQ memo. \\
\path{src/api/app.py} & Flask application factory and JSON API endpoints for borough summaries, boundaries, station footprints, and footprint scenarios. \\
\path{src/dashboard/dashboard.js} & D3 linked-view dashboard state, rendering, brushing, filtering, and table update logic. \\
\path{tests/} & Pytest contract tests for processed data, API behaviour, mobilisation matching, station footprints, station-address proximity, and report figures. \\
\bottomrule
\end{tabularx}
\end{center}

\subsection{Selected Implementation Listings}

Listings~\ref{lst:nullable-exceedance}--\ref{lst:dashboard-render-loop} show the parts of the implementation most directly tied to the thesis claims: denominator-safe response-time aggregation, bounded API filtering, and linked dashboard re-rendering. Full source files are submitted with the project code.

\begin{lstlisting}[language=Python, caption={Nullable exceedance flag used to keep missing first-pump times out of the denominator.}, label={lst:nullable-exceedance}]
rt = df["response_time_seconds"]
df["exceeds_six_min_target"] = pd.array(
    np.where(rt.isna(), pd.NA, rt > FIRST_PUMP_TARGET_SECONDS),
    dtype="boolean",
)
\end{lstlisting}

\begin{lstlisting}[language=Python, caption={Borough-summary API filtering rejects unknown keys before returning aggregate rows.}, label={lst:api-filtering}]
unknown = set(request.args.keys()) - set(FILTERABLE_COLUMNS)
if unknown:
    return jsonify({
        "error": "unknown filter keys",
        "unknown_keys": sorted(unknown),
        "allowed_keys": list(FILTERABLE_COLUMNS),
    }), 400

df = summary
for col in FILTERABLE_COLUMNS:
    values = request.args.getlist(col)
    if values:
        df = df[df[col].isin(values)]
\end{lstlisting}

\begin{lstlisting}[caption={Shared dashboard state drives the linked map, trend, distribution, and ranking panels.}, label={lst:dashboard-render-loop}]
function render() {
  renderSummary();
  renderMap();
  renderTrend();
  renderDistribution();
  renderTable();
}

path.on("click", (_, feature) => {
  const match = state.rows.find(
    (row) => normName(row.borough_canonical) === normName(feature.properties.name),
  );
  if (match) {
    state.borough = match.borough_canonical;
    els.borough.value = state.borough;
    render();
  }
});
\end{lstlisting}

\subsection{Reproducibility Commands}

The current outputs were validated with the following commands from the repository root:

\begin{verbatim}
python src/data/build_published_station_proximity.py
python src/data/build_report_evaluation_figures.py
python src/data/build_closure_cohort_context.py
python -m pytest tests/ -q
python -m compileall src
git diff --check
\end{verbatim}

The broader rebuild sequence is documented in \path{RUNBOOK.md}. Binary raw inputs and generated parquets are intentionally excluded from version control; tracked markdown/JSON sidecars provide the audit trail for schemas, data-quality checks, source URLs, and transformation notes.

\subsection{Verification Test Matrix}

The automated verification suite contains 79 pytest tests. The table groups them by the validation class they support; counts are taken from the current \path{tests/test_*.py} functions rather than estimated from prose.

\begin{center}
\small
\begin{tabularx}{\textwidth}{p{3.2cm}p{2.3cm}Xp{3.2cm}}
\toprule
Verification class & Test count & Representative invariants & Edge cases covered \\
\midrule
Canonical schema and published baseline bands & 14 & Required raw/derived columns, datetime parsing, response-time seconds/minutes consistency, row/date/borough counts, incident-group taxonomy, false-alarm and exceedance bands, precise/rounded coordinate bands. & Unit regression, coordinate-family regression, incident-group drift. \\
Borough-summary consistency & 11 & Unique borough/month/group keys, cell count within expected grid, incident-count sum against canonical table, percentile ordering, probability ranges, borough-set match, recorded-time denominator and Wilson-interval checks. & Empty-cell handling, denominator regression, non-monotone percentiles, small-$n$ and partial-month flags. \\
Mobilisation schema and matched-canonical partition & 18 & Raw/derived mobilisation columns, datetime and attendance-time units, first-arriving-pump logic, join share, first-pump time agreement with canonical incidents, borough/station population, \path{matches_canonical_incident} truth table. & Cross-boundary and unmatched rows kept out of matched-only deployment claims. \\
API contract and footprint scenario paths & 20 & Health/dashboard/boundary routes, borough-summary schema, filter round-trips, multi-value unions, JSON null safety, station-footprint and footprint-scenario schemas. & HTTP 400 on unknown filter keys, unknown station names, unknown scenario keys, and removing all stations; unknown filter values return empty HTTP 200. \\
Station-location and public-address invariants & 13 & Incident-derived station centroid count, London coordinate envelopes, non-negative IQRs, centroid-in-own-incident bounding boxes, public station row/geocode counts, known typo corrections, D1 comparison matching, borough proximity completeness. & Prevents footprint centroids being treated as building locations or routing-grade coverage. \\
Report figures and sidecars & 3 & Required report-side figures exist, are non-empty, have memo sidecars, and summary metrics keep key borough/coordinate/sensitivity values in expected bands. & Missing figure, stale memo, headline metric drift. \\
\bottomrule
\end{tabularx}
\end{center}

\begin{center}
\small
\begin{tabularx}{\textwidth}{p{3.1cm}Xp{3.2cm}}
\toprule
Scripted manual UI check & Expected observation & Record retained \\
\midrule
Linked selection & Selecting a borough or ranking row updates the map, trend, distribution, metrics, and table from the same state object. & Dashboard overview and component screenshots in \path{figures/dashboard/}. \\
Empty state & A deliberately unpopulated selection shows the amber empty-state banner instead of leaving broken or stale panels. & \path{fig_dashboard_empty_state.png}. \\
Filter round-trip & Borough, month, and incident-group controls preserve the selected state and render denominator-qualified values. & Strand 2 task script and screenshot-capture utility. \\
Responsive layout & The 390\,px mobile viewport has no horizontal overflow; desktop remains the primary comparative-map context. & \path{fig_dashboard_mobile_390.png}. \\
\bottomrule
\end{tabularx}
\end{center}

\section{Data Sources, Figures, and Evaluation Files}

\subsection{Data Source Index}

\begin{center}
\small
\begin{tabularx}{\textwidth}{p{4.0cm}p{3.6cm}X}
\toprule
Source & Local role & Notes \\
\midrule
LFB Incident Records \cite{lfb_incidents} & Primary incident table & 2024-01 to 2026-02 slice; canonicalised into \path{lfb_canonical_2024_2026.parquet}. \\
LFB Mobilisation Records \cite{lfb_mobilisations} & Appliance-deployment context & Used through the matched-canonical partition for deployment summaries. \\
ONS / GLA borough boundaries \cite{ons_boundaries} & Borough geometry & Used for borough aggregation and choropleth rendering. \\
LFB public station-address list \cite{lfb_stations} & Station-address proximity context & Historical postcode-geocoded coordinates; road travel times are unavailable. \\
postcodes.io \cite{postcodesio} & Public postcode geocoder & Used to convert station postcodes to approximate coordinates. \\
\bottomrule
\end{tabularx}
\end{center}

\subsection{Report Figure Index}

\begin{center}
\small
\begin{tabularx}{\textwidth}{p{5.2cm}X}
\toprule
Path & Contents \\
\midrule
\path{figures/dashboard/} (report folder) & Six dashboard screenshots: desktop overview, linked selection, trend panel, ranking table, empty-state banner, and mobile viewport. \\
\path{figures/early_context/} (report folder) & Front-loaded figures: motivation borough map and linked-view wireframe; the earlier coordinate-uncertainty diagram is superseded by the fields-by-analysis matrix in \S\ref{sec:background-granularity} (Table~\ref{tab:evidence-support}). \\
\path{figures/literature/} (report folder) & Three cited literature figures reproduced without alteration: PanViz~2.0 interface structure, MacEachren et al.'s uncertainty visual variables, and Taylor's three-panel closure-probability comparison. \\
\path{figures/results_evaluation/} (report folder) & Seven results/robustness figures generated from tracked artefacts, including the Taylor closure-cohort descriptive cross-check. \\
\path{figures/station_proximity/} (report folder) & Two D1 station-address proximity figures used in the evaluation chapter. \\
\bottomrule
\end{tabularx}
\end{center}

\subsection{Evaluation Files}

\begin{center}
\small
\begin{tabularx}{\textwidth}{p{5.4cm}X}
\toprule
Path & Contents \\
\midrule
\path{tests/evaluation/strand1_findings_summary.md} & Heuristic-walkthrough findings. \\
\path{tests/evaluation/strand2_thinkaloud_task_script.md} & Think-aloud task protocol for borough comparison and linked-view interpretation. \\
\path{tests/evaluation/strand3_rq5_uncertainty_probe.md} & Uncertainty and station-footprint/proximity probe. \\
\path{tests/evaluation/strand2_strand3_findings_summary.md} & Non-identifying compact think-aloud findings summary for Participants 1--3. \\
\path{tests/evaluation/observation_log_template.md} & Per-session observation template for task completion, panel choice, uncertainty noticing, and station-caveat interpretation. \\
\path{tests/evaluation/capture_dashboard_figures.py} & Screenshot-capture utility for the dashboard figures used in \S6. \\
\bottomrule
\end{tabularx}
\end{center}

\section{Dashboard User Guide}

This appendix is a short operating guide for a reader who wants to launch the dashboard and reproduce the inspection workflow evaluated in \S6, without reading the implementation chapter first.

\subsection{Requirements and Launch}

The dashboard requires Python 3.10+ and the dependencies in \path{requirements.txt}. From the repository root:

\begin{verbatim}
python -m venv .venv && source .venv/bin/activate
python -m pip install -r requirements.txt
python src/data/build_canonical.py
python src/data/build_borough_summary.py
flask --app src.api.app:create_app run --debug --port 5057
\end{verbatim}

The two build scripts regenerate the untracked parquet files from the raw LFB workbook under \path{data/}. The raw file is excluded because of repository size and source-refresh cadence, rather than an OGL v2 restriction; its URL and SHA-256 are recorded in the provenance sidecar. The dashboard is then served at \path{http://127.0.0.1:5057/dashboard}.

\subsection{Controls and Linked Behaviour}

The interface holds one shared selection state: \emph{borough $\times$ month $\times$ incident group}. It can be changed three ways: the three select controls at the top, clicking a borough on the map, or clicking a row in the ranking table. Any change re-renders all four panels together:

\begin{itemize}
    \item \textbf{Borough map (choropleth).} Fill colour encodes the author-derived share above 360 seconds for the selected month and incident group. Hover shows incident count, median, p90, recorded denominator, and 95\% Wilson interval.
    \item \textbf{Monthly trend panel.} Incident count and the response diagnostic across 25 complete months plus partial February 2026.
    \item \textbf{Response-distribution panel.} Median-to-p95 response-time ranges for the three incident groups in the selected cell, against the dashed six-minute reference line.
    \item \textbf{Ranking table.} Boroughs ordered by the response diagnostic, with recorded $n$, 95\% interval, and small-sample marker.
\end{itemize}

If the selected combination has no data, an amber empty-state banner replaces silent blank panels; adjusting any filter recovers a populated view.

\subsection{How to Read the Values}

The response proportion is computed over incidents with a recorded first-pump time. It is an author-derived diagnostic, not LFB's official incident failure rate. Read the estimate with recorded $n$, its Wilson interval, incident group, and partial-month marker. Precise-coordinate availability does not condition the borough result.

\subsection{Station Context and Its Limits}

Station context appears in the report and API rather than the dashboard. Assignment footprints locate historical responsibility centroids, while the postcode coordinates come from a historical address file. Neither source is a current estate register.

\subsection{Example Task}

To repeat the comparison task in Table~\ref{tab:task-evaluation-matrix}: select \emph{Fire} as the incident group, click \emph{Hillingdon} on the map, note its exceedance share and trend, then click an inner-London row such as \emph{Tower Hamlets} in the ranking table and compare the two boroughs' values, recorded denominators, confidence intervals, and incident counts.

\subsection{Known Limitations}

The dashboard has no one-click reset, direct demand ranking, year/weekday/hour filters, full response distribution, or station UI. These omissions are recorded as partial requirements in Table~\ref{tab:req-traceability}.
